%% 
%% Copyright 2019-2024 Elsevier Ltd
%% 
%% This file is part of the 'CAS Bundle'.
%% --------------------------------------
%% 
%% It may be distributed under the conditions of the LaTeX Project Public
%% License, either version 1.3c of this license or (at your option) any
%% later version.  The latest version of this license is in
%%    http://www.latex-project.org/lppl.txt
%% and version 1.3c or later is part of all distributions of LaTeX
%% version 1999/12/01 or later.
%% 
%% The list of all files belonging to the 'CAS Bundle' is
%% given in the file `manifest.txt'.
%% 
%% Template article for cas-sc documentclass for 
%% double column output.

\documentclass[a4paper,fleqn]{cas-sc}

% If the frontmatter runs over more than one page
% use the longmktitle option.

%\documentclass[a4paper,fleqn,longmktitle]{cas-sc}

%\usepackage[numbers]{natbib}
%\usepackage[authoryear]{natbib}
\usepackage[authoryear,longnamesfirst]{natbib}

%%%%%%%%%%%%%%%%%%%%%%%%%%%%%%%%%%%%%%%%%%%%%%%%%%%%
% Personal packages
\usepackage[version=4]{mhchem} % Formula subscripts using \ce{}
\usepackage{amssymb}
\usepackage{amsmath}
\usepackage{blindtext}
\usepackage{booktabs}
\usepackage{caption}

% Caption below table 
\usepackage{threeparttable}

% Long table and floating behavior handling
\usepackage{longtable}
% \usepackage{placeins}
\usepackage{hyperref}
% \hypersetup{breaklinks=true}

%%%%%%%%%%%%%%%%%%%%%%%%%%%%%%%%%%%%%%%%%%%%%%%%%%%%%

%%%Author macros
\def\tsc#1{\csdef{#1}{\textsc{\lowercase{#1}}\xspace}}
\tsc{WGM}
\tsc{QE}
%%%

% Uncomment and use as if needed
%\newtheorem{theorem}{Theorem}
%\newtheorem{lemma}[theorem]{Lemma}
%\newdefinition{rmk}{Remark}
%\newproof{pf}{Proof}
%\newproof{pot}{Proof of Theorem \ref{thm}}

\begin{document}
\let\WriteBookmarks\relax
\def\floatpagepagefraction{1}
\def\textpagefraction{.001}

% Short title
\shorttitle{}    

% Short author
\shortauthors{}  

% Main title of the paper
\title [mode = title]{Hydrogen-supply and feedstock sensitivity in pyrolysis-based sustainable aviation fuel production with South Korean context}  

% Title footnote mark
% eg: \tnotemark[1]
% \tnotemark[1] 

% Title footnote 1.
% eg: \tnotetext[1]{Title footnote text}
% \tnotetext[1]{} 

% First author
%
% Options: Use if required
% eg: \author[1,3]{Author Name}[type=editor,
%       style=chinese,
%       auid=000,
%       bioid=1,
%       prefix=Sir,
%       orcid=0000-0000-0000-0000,
%       facebook=<facebook id>,
%       twitter=<twitter id>,
%       linkedin=<linkedin id>,
%       gplus=<gplus id>]
% ---------------- Author 1 ----------------
\author[1]{Hyein Jung}[orcid=0009-0002-8027-3385]%[<options>]
\ead{976hannah@kist.re.kr}
\credit{Conceptualization; Methodology; Software; Formal analysis; Investigation; Writing - Original Draft; Visualization}

% ---------------- Author 2 ----------------
\author[1,2]{Kyeongsu Kim}[orcid=0000-0002-9383-7605]%[<options>]
\ead{kyeongsu@kist.re.kr}
\credit{Resources; Writing - Review \& Editing; Visualization; Funding acquisition}
\cormark[1]

% ---------------- Author 3 ----------------
\author[3,4]{Jong Min Lee}[orcid=0000-0001-5395-0067]%[<options>]
\ead{jongmin@snu.ac.kr}
\credit{Resources; Writing - Review \& Editing; Supervision; Funding acquisition}
\cormark[1]

% ---------------- Author 4 ----------------
\author[5]{Tae Hoon Oh}[orcid=0000-0003-1847-2164]%[<options>]
\ead{oh.taehoon@unist.ac.kr}
\credit{Conceptualization; Formal analysis; Writing - Review \& Editing; Supervision; Project administration}
\cormark[1]

% Corresponding author indication (add to whichever author applies)
% \cormark[1]

% ---------------- Affiliations ----------------
\affiliation[1]{organization={Clean Energy Research Center, Korea Institute of Science and Technology},
            addressline={5, Hwarang-ro 14-gil, Seongbuk-gu},
            city={Seoul},
            postcode={02792},
            country={Republic of Korea}}

\affiliation[2]{organization={Division of Energy and Environment Technology, KIST School, University of Science and Technology},
            addressline={5, Hwarang-ro 14-gil, Seongbuk-gu},
            city={Seoul},
            postcode={02792},
            country={Republic of Korea}}

\affiliation[3]{organization={Department of Chemical and Biological Engineering, Seoul National University},
            addressline={1, Gwanak-ro, Gwanak-gu},
            city={Seoul},
            postcode={08826},
            country={Republic of Korea}}

\affiliation[4]{organization={Institute of Chemical Process, Seoul National University},
            addressline={1, Gwanak-ro, Gwanak-gu},
            city={Seoul},
            postcode={08826},
            country={Republic of Korea}}

\affiliation[5]{organization={Graduate School of Carbon Neutrality, Ulsan National Institute of Science and Technology},
            addressline={50, UNIST-gil, Eonyang-eup, Ulju-gun},
            city={Ulsan},
            postcode={44919},
            country={Republic of Korea}}

% Footnote text
% \fntext[1]{}

% For a title note without a number/mark
%\nonumnote{}

% Here goes the abstract
\begin{abstract}
Sustainable aviation fuel (SAF) offers a near-term pathway for reducing aviation emissions because it can be deployed using existing aircraft and fuelling infrastructure without major modifications. Fast pyrolysis of lignocellulosic biomass is a promising alternative given its abundant supply, but its economic and environmental disadvantages hinder their industrial deployment.  
Its economic and environmental performance can be enhanced by balancing economies of scale, finite biomass availability, and decarbonization costs. In this study, the optimal production scale is identified, at which rising feedstock costs and single-train capacity limitations offset each other without depleting resources. 
On the other hand, the CORSIA greenhouse gas intensity threshold is satisfied by integrating green ammonia-derived hydrogen and post-combustion carbon capture and storage (CCS). At current stage, CCS is the most cost-effective decarbonization option, whereas green ammonia becomes economically competitive when the grey-hydrogen price exceeds a crossover range of \$1.59--2.91/kg H$_2$. This crossover reflects the strong influence of hydrogen supply costs on SAF production economics and provides a practical benchmark for evaluating the future competitiveness of green ammonia. 
Overall, the results show that the preferred decarbonization pathway depends on the required emission reduction level, as well as the relative costs of hydrogen and CCS. These findings establish a resource-constrained framework for SAF deployment in which plant scale, biomass availability, emission targets, and future hydrogen prices must be considered jointly.
\end{abstract}

% Use if graphical abstract is present
%\begin{graphicalabstract}
%\includegraphics{}
%\end{graphicalabstract}

% Research highlights
\begin{highlights}
\item Developed an integrated sustainable aviation fuel process combining biomass pyrolysis with green ammonia-derived hydrogen. 
\item Identified a minimum-efficient production scale by balancing economies of scale with finite biomass availability.
\item Demonstrated that green ammonia becomes economically favorable at a grey-hydrogen price of \$1.59--2.91/kg H$_2$.
\end{highlights}


% Keywords
% Each keyword is seperated by \sep
\begin{keywords}
Sustainable aviation fuel \sep Lignocellulosic biomass \sep Green ammonia \sep Techno-economic assessment \sep Life-cycle analysis 
\end{keywords}

\maketitle

% Main text
\section{Introduction}\label{sec:intro}
% Body of Section 1; \section{Introduction} is set in main-sc.tex.

% ¶1. Aviation needs a drop-in decarbonization route
Aviation decarbonization remains challenging, particularly for long-haul flights, which are difficult to electrify. Sustainable aviation fuel (SAF) offers a near-term pathway to reducing aviation emissions because it can be used in existing aircraft and fuelling infrastructure without requiring major modifications \citep{atsonios2015alternative, doliente2020bio}. This pathway is particularly relevant to the aviation indsutry's commitment to halve net CO$_2$ emissions by 2050 relative to 2005, following the adoption of this target by the International Air Transport Association (IATA). Under the Carbon Offsetting and Reduction Scheme for International Aviation (CORSIA), SAF must achieve life-cycle greenhouse gas emissions at least 10\% lower than those of conventional jet fuel to qualify for emissions reductions. 

% ¶2. Current SAF routes are limited; lignocellulosic pyrolysis is the alternative
Existing SAF production pathways face constraints associated with feedstock availability or process requirements. Hydroprocessed esters and fatty acids (HEFA) relies on waste oils and fats, which are projected to meet only about 70\% of SAF demand in 2030 \citep{doliente2020bio}. Gasification followed by Fischer--Tropsch synthesis, in contrast, requires high-temperature operation and syngas cleanup to remove contaminants \citep{jones2013process}. Fast pyrolysis of lignocellulosic biomass offers an alternative by utilizing forest and agricultural residues, which account for more than 60\% of global biomass potential \citep{wang2019biomass}. Moreover, lignin-derived compounds can contribute to the aromatic and cycloalkane fractions required in jet fuel \citep{son2022hybrid, escalante2022evaluation}.

% ¶3. Arm A: scale lowers conversion cost, but biomass supply pushes back
The economic viability of pyrolysis-based SAF production depends not only on conversion efficiency but also on the scale of the production facility and the cost of securing biomass feedstock. Increasing plant capacity can reduce unit capital costs as sourcing expands beyond nearby low-cost resources \citep{chireshe2025cost}. Previous bioenergy studies have demonstrated that the interaction between conversion economies of scale and biomass supply costs can determine the cost-optimal facility size \citep{de2017cost}. Accordingly, assuming a constant feedsotck price may overstate the economic benefits of plant expansion. This trade-off is particularly relevant to South Korea, where the economically usable supply of forest residues must be distinguished from the country's substantial forest area. Therefore, determining the cost-optimal scale of pyrolysis-based SAF production requires evaluating conversion costs against scale-dependent biomass supply costs under domestic resource constraints. 

% ¶4. Two levers with different cost structures: green ammonia and process CCS
Beyond plant scale, reducing the carbon intensity of pyrolysis-based SAF requires improvements in process efficiency and carbon management. Fast-pyrolysis bio-oil contains substantial oxygen and must undergo hydrodeoxygenation (HDO) and related upgrading steps to produce aviation-range hydrocarbons, adding hydrogen demand, energy consumption, and process emissions \cite{sorunmu2020review}. Therefore, process intensification through heat integration, improved conversion and product recovery, and post-combustion CO$_2$ capture and storage (CCS) can reduce the energy and emissions burdens of SAF production. However, the benefits of these measures depend on the process configuration and the extent of achievable emission reduction, motivating an integrated assessment of process design, energy use, and carbon capture rather than evaluating individual measures in isolation. 

The environmental performance of pyrolysis-based SAF also depends on the origin of the hydrogen required for bio-oil upgrading. Conventional steam methane reforming relies on natural gas and generates direct CO$_2$ emissions, while upstream methane emissions and incomplete capture can contribute to the life-cycle greenhouse gas footprint even when carbon capture is employed \citep{howarth2021green}. Green ammonia offers an alternative route for supplying renewable hydrogen, leveraging ammonia's potential for large-scale storage and transport while requiring additional energy for cracking and hydrogen recovery \citep{salmon2021green}. Comparing grey hydrogen with green-ammonia cracking under different post-combustion CO$_2$ capture configurations can thus establish how hydrogen sourcing and process-level carbon management jointly affect the economic and environmental performance of pyrolysis-based SAF. 

% ¶6. Objective and contributions
This study integrates techno-economic analysis (TEA) and life-cycle assessment (LCA) of a forest-residue fast-pyrolysis process for SAF to address two sequential questions: what plant capacity minimizes the minimum jet fuel selling price (MJSP) when biomass supply costs and equipment constraints are considered, and how should hydrogen supply and carbon capture be configured to meet the CORSIA emissions threshold at that capacity? The study makes three contributions. First, it identifies the cost-optimal plant capacity by accounting for single-train equipment limits and comparing three biomass pricing models: a constant price, a log-linear supply curve, and a fixed-elasticity supply curve. Second, at the selected capacity, it evaluates 30 configurations spanning the green-ammonia share of hydrogen supply and the post-combustion CO$_2$ capture rate, maps their costs and life-cycle emissions, identifies the lowest-cost configuration that meets the CORSIA threshold, and quantifies the cost of further emissions abatement. Third, it assesses the sensitivity of configuration rankings to price uncertainty using tornado and Shapley Additive exPlanations (SHAP) analyses, and identifies the grey-hydrogen price at which green-ammonia cracking becomes more cost-effective than CCS alone.

% ¶7. Paper organization
The remainder of this paper is organized as follows. Section~\ref{sec:methods} describes the process design, TEA, biomass feedstock pricing models, ammonia-derived hydrogen supply, CCS costing, and LCA. Section~\ref{sec:results} presents the optimal plant scale, the cost--emission trade-offs across hydrogen supply and CCS configurations, and the sensitivity and uncertainty analyses. Section~\ref{sec:conclusion} summarizes the main findings and conclusions.

\section{Methods}\label{sec:methods}
% Body of Section 2; \section{Methods} is set in main-sc.tex.

\subsection{Process design}\label{sec:methods-process}

\subsubsection{Main process unit}\label{sec:methods-main}

\begin{figure}[htbp]
    \centering
    \includegraphics[width=\textwidth]{BiomassLCA-BioPyrolysis (3).pdf}
    \caption{Biomass pyrolysis-to-SAF production process flow diagram. The diagram shows the major process units, biomass and hydrogen feed, intermediate streams, and product outputs in the analysis.}
    \label{fig:pfd-main}
\end{figure}

The process is based on the biomass fast-pyrolysis-to-fuels flowsheet developed by NREL/PNNL \citep{jones2013process}. The flowsheet comprises eight process areas: feed handling, pyrolysis, hydrotreating, product finishing, hydrogen recovery, hydrocracking, off-gas combustion, and the steam and balance-of-plant island (Fig.~\ref{fig:pfd-main}). Detailed description is available in SI~S1. 

In this study, the process design is adjusted from the early study of \citep{jones2013process} in four distinctive points. 
First, product slate of the process is shifted to maximize SAF production for the South Korean market, rather than the gasoline and diesel blendstock production targeted in the original design. Gasoline- and diesel-range streams are retained as value-generating co-products.
Second, surplus electricity is generated as a co-product via a gas-turbine power island, which is newly introduced in this study. 
Next, the original grey-hydrogen supply can be partially or fully replaced with hydrogen produced by green-ammonia cracking, as described in Section~\ref{sec:methods-cracker}.
Lastly, an amine-based CCS unit is added to capture recoverable CO$_2$ streams leaving combustion chambers or separators (Section~\ref{sec:methods-ccs}).

At the base scale, the plant processes 0.343~MMt/yr of woody biomass and produces 0.041~MMt/yr of SAF. The bio-oil sent to hydrotreating amounts to 0.244~MMt/yr, equivalent to approximately 71~wt\% of the biomass feed. Total fresh make-up hydrogen consumption is 0.024~MMt/yr, comprising 0.021~MMt/yr for hydrotreating and 0.003~MMt/yr for hydrocracking. Hydrotreating employs ruthenium catalysts in a series of two reactors at a liquid hourly space velocity (LHSV) of 0.5, followed by the last reactor with molybdenum catalyst at an LHSV of 0.22. Catalyst costs are estimated separately from reactor-vessel costs.

The Aspen Plus flow-sheet is costed using Aspen Process Economic Analyzer (APEA). The resulting \textit{Main} block, which represents whole-flow sheet capital and operating expenditures (CAPEX and OPEX) obtained from the APEA, is combined with five separately costed sub-units: CCS, pressure swing adsorption (PSA) hydrogen recovery, the off-gas reformer, hydrotreating catalysts, and the hydrocracker. Cost interpolation for these sub-units is detailed in SI~S3.

Plant capacity is scaled from the base-case SAF output of 0.041 MMt/yr to 1.23 MMt/yr, representing a 30-fold difference between the smallest and largest plant scales. Rather than applying a single cost-scaling exponent to the entire plant, each major unit is assigned an explicit single-train capacity limit. Additional parallel trains are introduced when a unit reaches its limit. The capacity limits and corresponding numbers of trains at each scale are reported in SI~S4. Because additional trains replicate existing units rather than enlarge individual units, they do not benefit from further single-unit economies of scale. This train-level configuration is implemented explicitly in the assessment at each plant scale.

All costs are reported in first-quarter 2022 USD. Costs from the original study of \citep{jones2013process} are escalated using a cost-index ratio of 3230/2250. The financial assumptions are 8,000 operating hours per year, a 30-year plant life, a 10\% internal rate of return (IRR), and a 35\% tax rate. The complete cost basis and financial assumptions are provided in SI~S2.

\subsubsection{Ammonia cracker}\label{sec:methods-cracker}

Green ammonia is used as a hydrogen carrier, with ammonia imported and catalytically cracked at the site to supply make-up hydrogen. The cracker capacity is determined by the plant's make-up hydrogen demand, with a single-train capacity limit of 403~GWh/yr. The reaction network, gas hourly space velocity (GHSV), and kinetics follow the source used to develop the Aspen model \citep{cho2024ammonia}. Ammonia conversion and hydrogen recovery also follow the parameters set in the original source. 

The recovery process was modified to employ a sulfuric acid scrubbing system to capture converted ammonia. It employs a multi-stage sulfuric acid scrubber to convert left ammonia into a liquid fertilizer \citep{manuzon2007prototype}. At the end, the scrubbed gas proceeds to nitrogen gas separation, yielding high-purity hydrogen (99.9\% H$_2$). The produced nitrogen fertilizer is credited as an ammonia-side co-product. Further details of the cracker process are provided in SI~S5.

\subsubsection{Carbon capture and storage}\label{sec:methods-ccs}

Post-combustion CCS is modeled using a monoethanolamine (MEA)-based retrofit. Its cost is interpolated from the power-plant amine-capture reference case reported by \citet{rao2002}, which processes 3.13~MMt/yr of CO$_2$, has a total CAPEX of \$297.9~M, requires 7,772~GWh/yr of regeneration heat, and consumes auxiliary electricity equivalent to 14\% of the heat duty. All costs are converted to refer to first-quarter 2022 USD.

The total captured mass comprises the exhaust-gas-recirculation (EGR) CO$_2$ stream and a specified fraction of the recoverable process CO$_2$ stream. The EGR CO$_2$ stream is completely captured. Capture-unit CAPEX scales with captured mass according to the six-tenths rule, whereas OPEX scales with an exponent of 0.7. Utility requirements scale linearly with captured mass. Capture OPEX includes make-up fuel for regeneration heat not supplied by on-site heat, auxiliary electricity, and optional CO$_2$ transport and storage (T\&S), the details for which is described in section~\ref{sec:methods-ccscost}). The capture fraction is varied from 0 to 1.0. The complete CCS cost derivation is provided in SI~S6.

\subsection{Techno-economic analysis}\label{sec:methods-tea}

The economic performance metric is the minimum jet-fuel selling price (MJSP,  \$/L SAF). MJSP is determined by discounted cash-flow analysis as the SAF selling price required to achieve a zero net present value given the economic parameters, assuming straight-line depreciation and zero salvage value.
Revenue from naphtha, diesel, exported electricity, and ammonia-side fertilizer is credited against the required revenue. Electricity and natural gas prices are set at \$65.0/MWh and \$49.1/MWh, respectively.

\subsubsection{Feedstock price model}\label{sec:methods-feed}

Because the domestic forest-residue supply is finite, the delivered biomass price is modeled as a function of plant scale rather than held constant. Larger plants are assumed to require progressively more expensive feedstock as regional supply is depleted and additional sources are accessed. 
% The regional supply context is described in SI~S8. 

% The scarcity of FR biomass feedstock was reflected in the scale-up analysis as a feedstock price via two main approaches. 
This study considers the full spectrum of woody biomass as viable feedstock, capitalizing on the proven versatility of pyrolysis to process both softwood and hardwood species. Literature consistently demonstrates successful bio-oil conversion across a diverse range of softwoods, such as pine and conifers \cite{tran2021fast, jarvis2014characterization}, as well as hardwoods, including black locust and birch \cite{straker2015black, li2014autothermal}. While comparative research indicates that compositional factors—specifically differences in lignin and ash content—may influence the precise yield profile \citep{garcia2007vacuum, azeez2010fast}, the fundamental feasibility of thermochemical conversion remains robust across species. Consequently, this process is designed to be feedstock-flexible, accommodating any woody biomass regardless of its classification.

The economic viability of this feedstock-flexible approach is further supported by the immense volume of available biomass. The total potential supply is substantial, comprising approximately 706~MMt of logs and 412~MMt of underutilized residues. However, these resources are subject to distinct market dynamics; currently, logs are valued at \$0.19/kg, while residues trade at a significantly lower price of \$0.11/kg. Given this vast scale and the sharp cost disparity between high-quality timber and lower-cost residues, static pricing data is insufficient for long-term strategic planning. Consequently, a robust price prediction model is required to accurately forecast market fluctuations and optimize feedstock procurement strategies.

The first model treats forest-residue biomass and log as interchangeable sources for products with wood fiber. 
This approach defines the supply curve using two key data points: the total available residues at a low price and the amount of residues and log supplies added at the current log price. 
The supply equation parameters are set to make the curve passes through both data points. 
However, this method often results in a supply curve that is more price-sensitive (\textit{elastic}) than what is typically observed in real markets, where pulpwood supply is usually less responsive to price changes. 

To address this discrepancy, the second approach uses a price elasticity adopted from previous research \cite{galik2009forest}.
This method assumes a constant, inelastic price response for pulpwood in the states of North Carolina, South Carolina, and Virginia with a fixed elasticity of 0.3 in a supply curve. 
It also assumes that all available forest-residue biomass is used up before any log is purchased for feedstock. 
The supply curve is constructed by following the same linear curve of the first model until the current log price and quantity. 
From then on, it maintains the elasticity of 0.3, making the slope steeper (SI~S7). 

\subsubsection{Green ammonia and hydrogen supply}\label{sec:methods-nh3}

Green ammonia is assigned a base price of \$635/t NH$_3$, with sensitivity analysis spanning \$270/t, representing a projected future price, to \$1,000/t, representing a current price. Hydrogen produced through ammonia cracking replaces grey hydrogen at a specified substitution fraction. The grey-hydrogen reference is priced at \$1,499/t H$_2$ \citep{katebah2022}.

\subsubsection{CCS cost}\label{sec:methods-ccscost}

Capture costs are calculated as described in Section~\ref{sec:methods-ccs}. T\&S costs are applied to the sequestered CO$_2$ mass at a base rate of \$10/t CO$_2$ , with a sensitivity range of \$4--45/t CO$_2$. The T\&S cost contribution varies linearly with its unit price.

\subsubsection{CO$_2$ abatement cost}
CO$_2$ abatement cost was used to evaluate the trade-off between hydrogen supply and CCS. Its definition is as follows: 
    \begin{equation}
    \text{Abatement cost} =
    \frac{\Delta\mathrm{MJSP}}
    {\Delta\mathrm{GHG}\times E_{saf}},
    \end{equation}
where \(E=31.66\)~MJ/L is the volumetric energy content of SAF. Because producers must at least meet the compliance threshold set by CORSIA, the reported value represents the additional cost of abatement beyond the reference configuration. Abatement costs values are undefined for configurations at or above the GHG intensity of compliance threshold. 


\subsection{Life-cycle assessment}\label{sec:methods-lca}

\subsubsection{System boundary and biogenic carbon}\label{sec:methods-boundary}

The functional unit is 1~MJ of SAF, and the system boundary covers the well-to-wake (WTWa) life cycle, comprising the well-to-pump (WTP) and pump-to-wake (PTWa) stages. The PTWa stage, representing fuel combustion, is adopted from the Fischer--Tropsch jet-fuel reference. Environmental impacts are expressed as GHG-100, based on 100-year global warming potential, in g~CO$_2$e/MJ SAF. Life-cycle inventories are derived from Greenhouse gases, Regulated Emissions, and Energy use in Technologies (GREET) model \cite{wang2023greet}, since it coincides with the standard of CORSIA. 

The total forest residue production emission includes forest maintenance activities, harvesting, collecting, bailing, and extracting, such as milling and chipping \cite{han2011well}.
In the meantime, ILUC emission of the feedstock is not considered since it has low risk in generating ILUC emissions, which applies similarly to biomass from wastes and microalgae \cite{ICAO2023CORSIA}.
This total emission of the production process is allocated to forest residue feedstock based on mass fractions, except for the emissions induced by fertilizer or agrochemical inputs \cite{han2011well}. 

Ammonia was assumed to be produced by renewable energy in Australia \cite{lee2022techno} and transported via conventional ship fueled by liquefied natural gas. Applying the technology of ammonia-fueled ships \cite{schwarzkopf2023future}, its benefit was negligible ($<$ 1.2 g CO$_2$e/MJ), so it was not included in the results of this paper.
Heat utilities such as mid-pressure and low-pressure steam were counted by heat consumption of fuel, which is natural gas from North America, without considering transportation. 
Heat exchange network in the process requires some make-up boiler feed water, and it is counted as process water resources. 

The counting of emissions from the downstream product chain imitates that of conventional synthetic paraffinic kerosene fuel. 
SAF produced from the plant is distributed via multiple routes following the assumption of the domestic jet distribution over the U.S. \cite{xie2011assessment}.
After reaching at pump stations, the PTWa emission profile is assumed to be the same as Fischer--Tropsch jet fuel from biomass \cite{elgowainy2012life}.
SAF compliance is assessed against the CORSIA requirement that eligible SAF achieve at least a 10\% reduction in WTWa GHG-100 intensity relative to conventional jet fuel.

Biogenic carbon is treated under the carbon-neutrality assumption. Biogenic combustion emissions are tracked separately as \(\text{CO$_{2, biogenic}$}\), distinct from fossil CO$_2$, such that biomass uptake and subsequent combustion emissions cancel. Biogenic CO$_2$ diverted to CCS is instead counted as a negative emission.

% \subsubsection{Allocation}\label{sec:methods-alloc}

Co-product emissions are allocated among the products for the central LCA estimate. This approach avoids strong dependence on displacement credits for exported electricity and fertilizer, whose results can vary substantially with the assumed displaced products and market conditions. 

% \subsubsection{SAF compliance criterion}\label{sec:methods-corsia}

\subsection{Sensitivity analysis}\label{sec:methods-sens}

% Three complementary analyses are conducted.

% First, cost-emission trade-off of hydrogen supply and CCS was evaluated.
% The main plant is fixed at the knee-point scale of 0.41~MMt/yr SAF. The green-ammonia hydrogen substitution fraction (0--1) and CCS capture fraction (0--1.0) are varied over a \(6\times5\) grid, yielding 30 configurations. Green NH$_3$ hydrogen shares of 0, 10, 25, 50, 75, and 100\% are crossed with CCS capture fractions of 0, 25, 50, 75, and 100\% (Table~\ref{tab:S9-4}). The non-dominated configurations are identified by simultaneously minimizing MJSP and well-to-wake GHG-100 intensity.

%%%% TO APPENDIX
% MJSP is calculated using the full TEA for every grid point. GHG intensity is evaluated on the allocation basis using the Korean electricity grid mix (Section~\ref{sec:methods-alloc}). It is calculated from the simulated CCS cases and a constant ammonia-related emissions adjustment:
%     \begin{equation}
%     \mathrm{GHG}(a,c) =
%     \mathrm{GHG}_{\mathrm{CCS}}(c)
%     -a\left(\mathrm{GHG}_{\mathrm{grey}}-\mathrm{GHG}_{\mathrm{NH}_3}\right),
%     \end{equation}
%     where \(a\) denotes the green-ammonia substitution fraction and \(c\) denotes the CCS capture fraction. This formulation assumes that the upstream emissions avoided by replacing grey H$_2$ are independent of, and additive to, the process CO$_2$ captured by CCS.

% Uncertainty bands for MJSP are constructed by varying the biomass price across the supply-curve range at this plant scale (107.3--116.1~\$/t), from the fixed forest-residue reference price to the supply-curve price (Section~\ref{sec:methods-feed}).

Sensitivity of process economics was compared by computing SHAP attribution of different parameters. Sensitivity analysis was applied only when the compliance threshold was satisfied based on LCA. For each configuration, 128 Latin hypercube samples are generated over the specified parameter ranges and evaluated using the TEA model. A gradient-boosting regression surrogate is fitted to the resulting samples, with an in-sample \(R^2>0.99\) required. TreeSHAP values are then calculated from the surrogate to attribute each sample's MJSP to individual drivers.

The varied parameters are grey-hydrogen price (\$1,000--4,000/t), green-ammonia price (\$270--1,000/t), CAPEX inflation (\(\pm30\%\)), CO$_2$ T\&S cost (\$4--45/t for CCS configurations), electricity price (\(\pm30\%\)), co-product prices (\(\pm20\%\)), and capacity factor (85\% or 91\%). Biomass price is additionally varied across the supply-curve range according to their production scale.

\begin{table}[htbp]
\centering
\caption{Varied parameters and corresponding ranges or values.}
\label{tab:varied_parameters}
\begin{tabular}{ll}
\hline
\textbf{Parameter} & \textbf{Range / Value} \\
\hline
Grey-H$_2$ price & \$1,000--4,000/t \\
Ammonia price & \$270--1,000/t \\
CAPEX & $\pm 30\%$ \\
CO$_2$ T\&S cost (CCS configurations) & \$4--45/t \\
Electricity price & $\pm 30\%$ \\
Co-product prices & $\pm 20\%$ \\
Capacity factor & 85\% or 91\% \\
Biomass price & Supply-curve range according to production scale \\
\hline
\end{tabular}
\end{table}


\section{Results and Discussion}\label{sec:results}

\subsection{Scale-up}
\label{sec:scale_feedstock}

Scale-up study in SAF production plant demonstrates the opposing effects of economies of scale and feedstock scarcity. Following the increase in process scale, economies of scale reduces CAPEX and overall conversion costs. On the other hand, biomass demand also increases, which raises the feedstock price as low-cost forest-residue resources become increasingly constrained. Their interaction determines the scale at which further expansion no longer reduces the MJSP.

\begin{figure}[htbp]
    \centering
    \includegraphics[width=\textwidth]{composite_scale_up_mjsp.png}
    \caption{Effect of SAF production scale and feedstock price on the MJSP. 
    (a) MJSP breakdown as a function of SAF production scale at a constant forest-residue price. 
    (b) Biomass price and MJSP under the forest-residue (FR) price, FR--Log Linear, and Fixed Elasticity models. 
    Bars represent biomass price, while lines represent the MJSP for the grey-hydrogen (H$_2$) base case, green-ammonia (NH$_3$) substitution, and CCS. Shaded regions indicate the range arising from the biomass price models.}
    \label{fig:scale_feedstock}
\end{figure}
%%%% Numbers should be updated based on the figure and SI

Increasing SAF production from 0.04 to 1.24~MMt/yr reduces the MJSP of the base process from \$1.65 to \$0.68/L SAF when the biomass price is fixed at the forest-residue floor (Fig.~\ref{fig:scale_feedstock}a; SI Table~S8-1). This reduction is primarily associated with the scale-sensitive components of the conversion cost. Because co-product revenue exceeds the raw-material cost at all considered scales, the remaining gross production cost is dominated by capital investment, utilities, and fixed operating costs. Capital investment accounts for approximately half of the total MJSP reduction up to 0.41~MMt/yr, while utilities and fixed operating costs account for most of the remainder. However, the scale benefit diminishes substantially with increasing capacity. Each capacity doubling reduces the MJSP by approximately 20\% below 0.41~MMt/yr, compared with only approximately 10\% above this scale. Consequently, 0.41~MMt/yr captures approximately 90\% of the total MJSP reduction obtained by increasing the capacity to 1.24~MMt/yr. The diminishing scale benefit is also associated with the increasing capital scaling exponent, which approaches unity as the pyrolyzer, gas turbines, compressors, and expanders reach single-train capacity limits and require parallel trains (SI~S3-4).

This scale effect changes when the increase in biomass demand is coupled with a realistic feedstock supply curve (Fig.~\ref{fig:scale_feedstock}b; specific numbers in SI Table~S8-2). Below approximately 0.41~MMt/yr, all three biomass-price models remain at the forest-residue price floor of \$0.11/kg because the biomass demand remains within the low-cost forest-residue supply. At the production scale of 0.41~MMt/yr, the biomass price increases only to \$0.12/kg, corresponding to an 8.3\% increase and an MJSP increase from \$0.79 to \$0.85/L for the base process. At larger scales, however, the feedstock-price increase progressively offsets the scale benefit.  

Therefore, the opposing effects of scale and feedstock price produce a minimum MJSP at 0.41~MMt/yr scale under both rising feedstock-price models. Below this scale, the reduction in conversion cost from increasing capacity exceeds the increase in feedstock cost; above this scale, the feedstock-price penalty becomes larger than the additional scale benefit. Therefore, increasing the capacity beyond this range provides limited economic benefit while exposing the process to increasing feedstock costs. 

At this critical scale, hydrogen supply and process configuration become important considerations. A configuration change is essential, because the base configuration using grey hydrogen remains above the CORSIA emission threshold. However, the green-ammonia and CCS configurations consistently incur higher production costs across the considered scales and feedstock price models. Notably, at a scale of 0.41~MMt/yr, operating expenditure accounts for more than 80\% of the gross production cost, indicating that hydrogen supply and process configuration offer greater opportunities for cost reduction than further plant-scale expansion.

\subsection{Cost--emission trade-off of hydrogen supply and CCS}
\label{sec:h2_ccs_tradeoff}

The CORSIA compliance threshold can be achieved through post-combustion CCS alone, with a minimum capture rate of 60\% (Fig.~\ref{fig:pareto_crossover_shap}a). The CCS 60\% configuration achieves a WTWa GHG intensity of 75.14~g~CO$_2$e/MJ at an MJSP of \$0.977/L, making it the lowest-cost compliant configuration among the evaluated options. This result highlights the economic advantage of CCS at moderate capture rates, which provides substantial emission reductions at relatively low incremental cost.

In contrast, green-ammonia substitution alone is economically unfavorable for achieving compliance. Each 10-percentage-point increase in green-ammonia substitution raises the MJSP by \$0.185/L, corresponding to an incremental abatement cost of \$1,530/t~CO$_2$e avoided (Fig.~\ref{fig:pareto_crossover_shap}b). These results indicate that green-ammonia substitution alone is both insufficient for compliance and costly relative to CCS.

Combining green-ammonia substitution with CCS provides an alternative pathway to compliance. For example, a configuration with 10\% green-ammonia substitution and 50\% CCS achieves a comparable GHG intensity to the CORSIA threshold, with an MJSP of \$1.098/L compared with \$0.977/L for CCS alone. Although this combination does not offer a cost advantage under the current assumptions, it demonstrates that compliance can be achieved through different combinations of hydrogen supply and carbon capture. Therefore, the preferred configuration depends on the relative costs of green ammonia and CCS, as well as the desired degree of emission reduction beyond the compliance threshold.

% The plant capacity was subsequently fixed at 0.42~MMt/yr to examine the effects of hydrogen source and CCS. Green-NH$_3$ hydrogen substitution and CCS generate distinct cost--emission trajectories, spanning an MJSP of \$0.79--3.09/L and a WTWa GHG intensity of 100.8--21.9~g~CO$_2$e/MJ SAF (Fig.~\ref{fig:pareto_crossover_shap}a). Thus, the lowest-emission configuration reduces GHG intensity by approximately 4.6-fold relative to the grey-H$_2$ base case, but increases the MJSP by approximately 3.9-fold.

\begin{figure}[htbp]
    \centering
    \includegraphics[width=\textwidth]{composite_pareto_crossover_shap.png}
    \caption{Cost, emissions, and MJSP sensitivity to hydrogen supply and CCS. 
    (a) MJSP versus WTWa GHG intensity for combinations of green-ammonia (NH$_3$) hydrogen share and CCS capture fraction. 
    (b) CO$_2$ abatement cost relative to the CCS 60\% compliant base over the same green-ammonia (NH$_3$) hydrogen share and CCS capture fraction grid; blank cells emit at or above the base and have no
  abatement cost.
    (c,d) SHAP-based attribution of MJSP drivers for CCS 60\% and Green-ammonia 10\% + CCS 50\%, respectively. 
    (e) MJSP of the two compliant configurations as a function of grey-hydrogen price at different green-ammonia prices. 
    The dashed line in (a) and (e) denotes the CORSIA threshold and base grey-hydrogen price, respectively.}
    \label{fig:pareto_crossover_shap}
\end{figure}

% Green-NH$_3$ substitution produces an approximately linear change in both cost and emissions. Increasing the green-NH$_3$ share by 10 percentage points decreases the WTWa GHG intensity by 3.81~g~CO$_2$e/MJ while increasing the MJSP by \$0.185/L, independent of the CCS level. This corresponds to an incremental cost of \$1,530/t~CO$_2$e avoided, resulting in approximately parallel iso-CCS trajectories in the cost--emission space. 

% In contrast, the CCS effect is nonlinear, with a distinct change in the incremental cost after 50\% capture. Increasing post-combustion carbon capture rate from 0 to 50\% reduces the GHG intensity by 21.8~g~CO$_2$e/MJ for an MJSP increase of only \$0.12/L, corresponding to \$174/t~CO$_2$e avoided. Increasing capture from 50 to 100\% provides a further reduction of 19.0~g~CO$_2$e/MJ but requires an additional \$0.42/L, increasing the incremental cost to \$698/t~CO$_2$e avoided. Thus, the marginal cost of CCS increases substantially beyond 50\% capture, whereas the marginal cost of green-NH$_3$ substitution remains constant.

% The different cost trajectories result in distinct pathways to CORSIA compliance. The CORSIA threshold is 75.93~g~CO$_2$e/MJ, corresponding to 90\% of the WTWa GHG intensity of conventional jet fuel. The grey-H$_2$ base case has a GHG intensity of 100.8~g~CO$_2$e/MJ and therefore does not meet the threshold. CCS alone reaches the threshold at approximately 58\% capture. The CCS 60\% configuration has a GHG intensity of 75.14~g~CO$_2$e/MJ and an MJSP of \$0.977/L, making it the lowest-cost compliant configuration among the considered composites of green-NH$_3$ and CCS.

% On the similar GHG emission rate level, addition of green-NH$_3$ 10\% to CCS 50\% rate configuration also exists with MJSP at \$1.098/L. This is comparable to the lowest-cost compliant configuration using CCS only, from both of economic and environmental perspective.

% Beyond the compliance point, the incremental cost of further GHG reduction depends strongly on the selected lever(Fig.~\ref{fig:pareto_crossover_shap}b). Starting from the CCS 60\% configuration, increasing CCS to 75 and 100\% requires approximately \$733 and \$742/t~CO$_2$ avoided, respectively. In comparison, adding green-NH$_3$ results in substantially higher incremental costs of approximately \$900--1360/t~CO$_2$ over the 10--100\% green-NH$_3$ range along the CCS 75 and 100\% pathways. The full green-NH$_3$ and CCS configuration reaches the lowest simulated GHG intensity of 21.9~g~CO$_2$e/MJ at an MJSP of \$3.09/L, corresponding to approximately \$1250/t~CO$_2$ avoided relative to the CCS 60\% reference.

\subsection{Sensitivity of the hydrogen-supply and CCS trade-off}
\label{sec:sensitivity_tradeoff}

% The uncertainty in the cost--emission trade-off is dominated by hydrogen-related prices. Applying the parameter ranges defined in Section~2.2.4 gives MJSP ranges of \$0.55--1.84/L for the grey-H$_2$ base case and \$1.91--4.26/L for the full green-NH$_3$ + CCS 100\% configuration. Under favourable hydrogen prices and high capacity factor, the overall trade-off shifts downward by approximately \$0.3/L while retaining a similar shape.

SHAP analysis was performed for the two compliant configurations with the lowest MJSP: CCS 60\% and Green-ammonia 10\% + CCS 50\% (Fig.~\ref{fig:pareto_crossover_shap}c,d; SI Table~S8-6). The biomass contribution remains comparatively narrow because the biomass-price range at the selected plant scale is only \$0.11--0.12/kg. 
Drivers are ranked by their mean absolute SHAP values in Fig.~\ref{fig:pareto_crossover_shap}c,d. One-at-a-time MJSP ranges over the same parameter intervals are also calculated in SI (Table~\ref{tab:S9-7}).

Grey-hydrogen price was the dominant MJSP driver in both configurations, with mean absolute SHAP values of \$0.310/L for CCS 60\% and \$0.274/L for Green-ammonia 10\% + CCS 50\%. These values are approximately three times larger than those of the next most influential variables. Higher grey-hydrogen prices consistently increase the MJSP, producing the broad horizontal spread observed in both SHAP distributions. The larger influence in the CCS-only configuration reflects its complete dependence on grey hydrogen, compared with 90\% grey-hydrogen dependence in the green-ammonia configuration.

CAPEX, co-product price, and electricity were the next most influential variables, with mean absolute SHAP values of approximately \$0.10/L, \$0.08--0.10/L, and \$0.05/L, respectively. Capacity factor and biomass price had substantially smaller effects, quite negligible. The route-specific variables were also secondary. CO$_2$ T\&S cost parameter contributed approximately \$0.025/L and \$0.023/L for CCS 60\% and Green-ammonia 10\% + CCS 50\%, respectively. Green-ammonia price contributed \$0.031/L in the latter configuration because green ammonia supplied only 10\% of the total hydrogen demand. Thus, although the choice of compliance route changes the contributions of route-specific variables, hydrogen price remains the primary source of MJSP uncertainty for both configurations.


\subsection{Green-ammonia--CCS cost crossover}
\label{sec:h2_crossover}

The resulting MJSP ranges of neighboring configurations overlap substantially, when the relative cost of the two compliance routes changes according to the varying hydrogen and CO$_2$ management prices (Fig.~\ref{fig:pareto_crossover_shap}e).
This indicates that the relative cost of different hydrogen-supply and CCS configurations depends on the assumed market conditions.

In the meantime, the two compliance configurations exhibit different sensitivities to grey-hydrogen price because of their different hydrogen compositions. The CCS-only configuration uses 100\% grey hydrogen and its MJSP increases by \$0.419/L for each \$1/kg increase in grey-hydrogen price. The Green-ammonia 10\% + CCS 50\% configuration uses 90\% grey hydrogen and therefore has a lower slope of \$0.378/L per \$1/kg~H$_2$. Consequently, increasing grey-hydrogen price progressively narrows the cost difference between the two configurations.

The crossover grey-hydrogen price depends strongly on the green-ammonia price. Only the two lowest green-ammonia price cases cross within the investigated grey-hydrogen range--crossover occurs at \$1.59~/kg~hydrogen for \$270/t NH$_3$ and \$2.91~/kg~H$_2$ for \$441/t NH$_3$ (SI Table~S8-5). At the base prices of \$1.50~/kg~H$_2$ for grey hydrogen and \$635/t for green ammonia, the CCS-only configuration remains lower in MJSP throughout the investigated grey-hydrogen range. The green-ammonia pathway becomes cost-competitive when the grey-hydrogen price increases sufficiently or when the green-ammonia price decreases toward approximately \$450/t, at which point the crossover occurs below approximately \$3~/kg~H$_2$.
This is substantially larger effect than CO$_2$ T\&S cost result of which is handled in SI Table~S8-7. 

% Across the considered green-NH$_3$ price range of \$270--1000/t, the one-at-a-time MJSP variation is \$0.234/L, compared with \$0.165--0.187/L for the CO$_2$ transport and storage range of \$4--45/t~CO$_2$. The resulting crossover in grey-H$_2$ price shifts by approximately 5.6~\$/kg~H$_2$ across the green-NH$_3$ price range, whereas the full CO$_2$ transport and storage range changes the crossover by no more than approximately 0.45~\$/kg~H$_2$ (SI Table~S9-7).


\section{Conclusion}\label{sec:conclusion}
%% Conclusion

This study demonstrates that the economic feasibility and decarbonization of forest-residue-derived SAF depend on jointly optimizing plant scale, biomass availability, hydrogen supply and carbon capture strategies. Scale-up from 0.04 to 1.24 MMt/yr SAF production rate reduces MJSP by approximately 60\%, bringing it below \$1/L. However, approximately 90\% of the total cost reduction is achieved at 0.41 MMt/yr production scale, where economies of scale and the increasing cost of biomass supply intersect. Beyond this scale, single-train capacity limitations and feedstock price escalation offset further economies of scale, highlighting the need to balance market substitution potential against domestic forest-residue availability.

The grey-hydrogen base case exhibits GHG intensity of 100.8 g CO$_2$e/MJ, exceeding the CORSIA threshold of 75.93 g CO$_2$e/MJ. Integrating green ammonia-derived hydrogen and post-combustion CCS reduces emissions to as low as 21.9 g CO$_2$e/MJ, although the MJSP increases from \$0.79/L to \$3.09/L across the evaluated design space. CCS provides the more cost-effective decarbonization pathway per unit of GHG emissions avoided, making the 60\% CCS configuration the lowest-cost option for achieving CORSIA compliance. 

Green ammonia-derived hydrogen becomes economically competitive with CCS when the grey-hydrogen price exceeds a crossover point of \$1.59-2.91/kg H$_2$. This behavior reflects the strong influence of hydrogen supply costs on the MJSP. The economic advantage of green ammonia therefore depends primarily on its cost relative to grey hydrogen. The crossover price provides a simple benchmark for evaluating when green ammonia becomes more competitive than CCS.

This relationship provides a transparent basis for evaluating how future changes in hydrogen prices may alter the preferred decarbonization strategy. Overall, the findings establish a scale-dependent and cost-constrained framework for SAF deployment that integrates finite biomass resources with alternative decarbonization pathways. The results underscore that neither maximum plant scale nor maximum emission reduction alone defines the preferred configuration; instead, deployment decisions must jointly account for feedstock availability, production economics, and emissions compliance.



% Numbered list
% Use the style of numbering in square brackets.
% If nothing is used, default style will be taken.
%\begin{enumerate}[a)]
%\item 
%\item 
%\item 
%\end{enumerate}  

% Unnumbered list
%\begin{itemize}
%\item 
%\item 
%\item 
%\end{itemize}  

% Description list
%\begin{description}
%\item[]
%\item[] 
%\item[] 
%\end{description}  

% Uncomment and use as the case may be
%\begin{theorem} 
%\end{theorem}

% Uncomment and use as the case may be
%\begin{lemma} 
%\end{lemma}

%% The Appendices part is started with the command \appendix;
%% appendix sections are then done as normal sections
% \appendix
% 
\documentclass[11pt]{article}

\usepackage[margin=1in]{geometry}
\usepackage{graphicx}
\usepackage{amsmath}
\usepackage{booktabs}
\usepackage{siunitx}

%%%%%%%%%%%%%%%%%%%%
% Personal packages
\usepackage{amssymb}
\usepackage{blindtext}
\usepackage{caption}

% Caption below table 
\usepackage{threeparttable}

% Long table and floating behavior handling
\usepackage{longtable}
% \usepackage{placeins}
\usepackage{xurl}
\usepackage{hyperref}

\usepackage[authoryear,longnamesfirst]{natbib}

%%%%%%%%%%%%%%%%%%%%

\title{
    Hydrogen-supply and feedstock sensitivity in pyrolysis-based sustainable aviation fuel production with South Korean context\\[1em]
    \Large\textbf{Supporting Information}
}
\author{Hyein Jung, Kyeongsu Kim, Jong Min Lee, Tae Hoon Oh}
\date{}

\begin{document}

\maketitle

\renewcommand{\thesection}{S\arabic{section}}
\counterwithin{table}{section}
\renewcommand{\thetable}{\thesection-\arabic{table}}

\section{Process Design Details}\label{sec:S0}

The pyrolysis process employs a circulating fluidized bed (CFB) reactor, selected for its scalability and proven performance in industrial applications by Ensyn \cite{ensyn_quebec} and Fortum \cite{fortum_biooil_joensuu}. 
The CFB reactor operates at 925~$^{\circ}$F, 21 psi with residence time in seconds to maximize bio-oil yield up to 75 wt\% while minimizing char formation.
Another key design parameter is the heat transfer regime, which utilizes sand as the fluidizing medium in this work. 
The sand enters the pyrolyzer with a temperature of 1,190~$^{\circ}$F and a sand-to-biomass ratio of 14.5 to ensure rapid heating rates and suppress secondary reactions. 

The pyrolysis product consists of three phases: crude bio-oil, solid char, and non-condensable gas (NCG). 
Crude bio-oil is a complex mixture of over 300 oxygenated compounds, modeled using representative groups from \cite{jones2013process}.
The model compounds used to represent each group are summarized in Table \ref{tab:model-compounds}. 
It is separated via a high-efficiency cyclone with 99.9\% capture rate, where heat-transferring sand is also separated.
Solid stream with char and sand is fed to the combustor, where the sand is heated up to 1189~$^{\circ}$F, ready to be sent back to the pyrolyzer.
NCG comprises CO, CO\textsubscript{2}, and light hydrocarbons as shown in Table \ref{tab:pyro-prod}.
A bi-level quenching system rapidly cools the vapor stream from 800~$^{\circ}$F to 113~$^{\circ}$F to stabilize reactive intermediates and remove NCG products from the bio-oil stream. 
The separated NCG components are compressed and fed back to the pyrolyzer as fluidizing gas. 

\begin{table}[htbp]
\centering
\caption{Composition of pyrolysis products, gas components, and bio-oil components.}
\begin{tabular}{lc||lc}
\toprule
\multicolumn{2}{c||}{\textbf{Pyrolysis Product}} & \multicolumn{2}{c}{\textbf{}} \\
\multicolumn{2}{c||}{(per dry biomass)} & \multicolumn{2}{c}{\textbf{Gas}} \\
\hline
\textbf{Components} & \textbf{wt\%} & \textbf{Components} & \textbf{wt\%}\\
\hline
Bio-oil & 64 & CO & 46 \\
Reaction water & 12 & CO\textsubscript{2} & 43 \\
Char \& ash & 12 & CH\textsubscript{4} & 6 \\
Gas & 12 & C\textsubscript{2}+ & 4.7 \\
& & H\textsubscript{2} & 0.3 \\
\midrule
\multicolumn{4}{c}{\textbf{Bio-oil Components}} \\
\hline
\textbf{Groups} & \textbf{wt\%} & \textbf{Groups} & \textbf{wt\%}\\
\hline
Acids & 5.57 & LMW Lignin A & 10.95 \\
Alcohols & 6.97 & LMW Lignin B & 2.21 \\
Ketones & 8.29 & Extractives & 2.99 \\
Aldehydes & 6.97 & HMW Lignin A & 9.15 \\
Guaiacols & 4.98 & HMW Lignin B & 1.99 \\
LMW Sugars & 5.97 & N-compounds & 0.07 \\
HMW Sugars & 33.86 & S-compounds & 0.025 \\
\bottomrule
\end{tabular}
\label{tab:pyro-prod}
\end{table}

Crude bio-oil requires upgrading to reduce its high oxygen content at 35~wt\% and acidity for fuel compatibility \cite{elliott2007historical}, achieved through a three-stage hydrotreating system \cite{elliott2016transportation}. First, a stabilizer operating at 140°C and 1200 psi with Ruthenium sulfide catalyst selectively hydrogenates reactive carbonyls and dienes, consuming minimal hydrogen (0.08 wt\%) to produce intermediates with lower acidity. The first-stage hydrotreater then processes moderate-reactivity oxygenates at 260°C and 2065 psi using Ru-based catalysts, requiring 7 wt\% hydrogen for temperature control via quenching and 0.18 wt\% for deoxygenation. Finally, sulfided molybdenum-based catalysts in the last stage (352°C, 2000 psi) saturate stubborn oxygenates and aromatics into cyclohexanes and iso-paraffins. Overall, the process converts 28.5 wt\% of dry biomass into upgraded oil with 12.2 wt\% hydrogen consumption. The composition of the upgraded oil was modeled  using user-defined components, based on experimental data reported in a previous study \cite{jones2013process}. Thermodynamic description of the user-defined components used in this study is in Table \ref{tab:app-user-define1}-\ref{tab:app-user-define2}.

After the sequential hydrotreating, the light column (T-290) separates NCG components smaller than C\textsubscript{5}.
In the design process, it separates pentane (light key component) from short-chain paraffins (heavy key components). 
The column employs a partial vapor condenser operating at 650 psi, which incorporates water entrainment to mitigate fouling and enhance separation performance. 
More detailed specification follows the description in Table \ref{tab:distil}.

\begin{table}[tbp]
\centering
\caption{Specifications for product distillation columns} \label{tab:distil}
\begin{threeparttable}
\begin{tabular}{l|c|c|c|c}
\toprule
\textbf{} & \textbf{T-290} & \textbf{T-310} & \textbf{T-320} & \textbf{T-330} \\
\midrule \midrule
Distillate & $<$ C\textsubscript{5} & Gasoline & SAF & Diesel \\
% \hline
% Light Key & Cyclo C6-HCK product & Cyclo C6 B & 1,2,3,4-Naphthalenes \\
% \hline
% Heavy Key & Cyclo C5 B & Phenols & Extractive intermediates \\
% \hline
No. of Stages & 30 & 30 & 20 & 20 \\
Feed stage & 8 & 15 & 10 & 10 \\
Reflux ratio\tnote{a} & 0.952 & 3.6 & 0.43 & 0.72 \\
Distillate/Feed\tnote{a} & 0.2 & 0.33 & 0.60 & 0.56 \\
Bottom Temp. & 367~$^{\circ}$F & 422~$^{\circ}$F & 512~$^{\circ}$F & 672~$^{\circ}$F \\
Condenser pressure & 65~psi & 45~psi & 20~psi & 15~psi \\
\bottomrule
\end{tabular}
\begin{tablenotes}
\footnotesize
\item[a] based on molar flow rate
\end{tablenotes}
\end{threeparttable}
\end{table}

    This study proposes sequential fractionation of hydrotreated bio-oil into three kinds of liquid fuels--gasoline, diesel, and SAF.
    % Upgraded bio-oil is fractionated into three kinds of liquid fuels - gasoline, diesel, and SAF. 
    The first column (T-310) for gasoline-range fuel operates at 45 psi, with cyclo C6-HCK products group as light key components and cyclo C5-B group as heavy components. 
    These hydrotreated compounds are defined in Table.\ref{tab:compounds-from-hydroprocess}.
    The last column (T-330) produces diesel-range fuel, separating 1,2,3,4-naphthalenes group from extractive intermediates.  
    The leftover long-chain hydrocarbons pass hydrocracker to produce short-chain products such as cyclo C6-HCK products group and return to the column T-310. 

%% Simulation: Pyrolysis_Jones_551_into_571_forQuality_TSOL2_retune.apw
\begin{table}[tb]
\centering
\caption{Comparison of Jet A and pyrolysis-based SAF product properties.}
\begin{threeparttable}
\begin{tabular}{l r r}
\toprule
\textbf{Property} & \textbf{Jet A} & \textbf{Pyrolysis SAF} \\
\midrule
\multicolumn{3}{l}{\textbf{Performance requirements}\tnote{a}} \\
% \midrule
Heat of combustion (MJ/kg) & $> 42.8 $ & 44.8 \\
Specific gravity  & 775--840  & 786 \\
Freezing point ($^\circ$C) & $< -40 $ & $-$56 \\
% Flash point, max ($^\circ$C) & $> 38 $ & 27.8 \\
Aromatic compounds (vol\%) & 8--25  & 10.2  \\
Sulfur content (wt\%) & $< 0.3 $ & NA \\
\midrule
\multicolumn{3}{l}{\textbf{Atomic composition comparison}\tnote{b}} \\
% \midrule
Carbon  (wt\%) & 86.4 & 83.8 \\
Hydrogen  (wt\%) & 13.5  & 13.6  \\
Oxygen  (wt\%) & 0.0044 & 0.026  \\
\bottomrule
\end{tabular}
\begin{tablenotes}
\footnotesize
\item[a] Ref. \cite{lahijani2022upgrading} 
\item[b] Ref. \cite{zhang2015production}
\end{tablenotes}
\end{threeparttable}
\label{tab:saf-spec}
\end{table}

    The second column (T-320) produces SAF within the C\textsubscript{8-16} hydrocarbon range, which is targeted to meet the performance specification of Jet A fuel \cite{Davidson2014}. 
    The evaluation of fuel quality was based on several key performance indicators: heat of combustion, flash point, specific gravity, and aromatic content \cite{lahijani2022upgrading}, as outlined in Table \ref{tab:saf-spec}.
    The heat of combustion was assessed conservatively using the lower heating value of the fuel components, while the flash point was determined by the lowest boiling point of the constituent compounds. 
    Specific gravity and volumetric fraction of aromatic compounds were evaluated via Aspen Plus simulations.
    The freezing point was estimated according to carbon class composition, employing correlation models proposed by \cite{shi2017quantitative}.
    Among these models, coefficients optimized by a genetic algorithm were selected, as they demonstrated the highest accuracy for freezing point estimation in catalytic fast pyrolysis-derived SAF \cite{watanasiri2023prediction}. 

    With the exception of the volumetric fraction of aromatic compounds, all evaluated properties were found to be within the desired range. 
    To address this, phenols were separated from the final product, which resulted in an aromatic content of 10.2 vol\%.
    Additionally, the atomic composition of carbon, hydrogen, and oxygen in the produced fuel was compared to that of conventional jet fuel \cite{zhang2015production}, confirming that the composition of T-320 product is comparable to standard aviation fuels.
    
    The off-gas collected from hydroprocesses, hydrotreating, and hydrocracking is rich in hydrogen gas.
    Hence, the leftover hydrogen gas is recovered via a pressure swing adsorption (PSA) unit to satisfy the high purity ($> 99$ wt\%) required for the hydroprocesses.
    Hydrogen recovery is assumed to be 85\%, referring to the industrial technology providing approximately 80\% recovery rate these days \cite{mahler_hydroswing}.
    The recovered hydrogen gas is recycled for the hydroprocesses with makeup feed from external sources. 
    
    The off-gas consists of 74.8 vol\% hydrogen gas, so the state-of-the-art gas turbine technology should be applied to resolve the high flame temperature and flame speed of hydrogen gas combustion.
    Exhaust gas recirculation (EGR) framework was applied to fulfill the maximum combustion temperature of 3002~$^{\circ}$F and prevent NOx emission \cite{dennis2024literature}.
    The gas turbine generates 65 MW with high-temperature waste heat, which is utilized in the distillation column reboilers, EGR preheating, and CCS absorber regeneration. 

Table \ref{tab:model-compounds} provides the information of bio-oil model compounds used in process modeling. 

\begin{table}[htbp]
\centering
\caption{Model compounds and their composition on a dry weight basis retrieved from \cite{jones2013process}}
\begin{tabular}{|p{4.5cm}|p{4.5cm}|l|r|}
\hline
\textbf{Compound group} & \textbf{Model compound} & \textbf{Formula} & \textbf{\% wt} \\
% \textbf{} & \textbf{} & \textbf{} & \textbf{(dry basis)} \\
\hline
\multicolumn{4}{|l|}{\textbf{\textit{Water soluble}}} \\
\hline
Acids & Crotonic acid & C\textsubscript{4}H\textsubscript{6}O\textsubscript{2} & 5.57 \\
\hline
Alcohols & 1,4-Benzenediol & C\textsubscript{6}H\textsubscript{6}O\textsubscript{2} & 6.97 \\
\hline
Ketones & Hydroxyacetone & C\textsubscript{3}H\textsubscript{6}O\textsubscript{2} & 8.29 \\
\hline
Aldehydes & 3-Methoxy-4-Hydroxybenzaldehyde & C\textsubscript{8}H\textsubscript{8}O\textsubscript{3} & 6.97 \\
\hline
Guaiacols & Isoeugenol & C\textsubscript{10}H\textsubscript{12}O\textsubscript{2} & 4.98 \\
\hline
Low MW sugars & Levoglucosan & C\textsubscript{6}H\textsubscript{10}O\textsubscript{5} & 5.97 \\
\hline
High MW sugars & Cellobiose & C\textsubscript{12}H\textsubscript{22}O\textsubscript{11} & 33.86 \\
\hline
\multicolumn{4}{|l|}{\textbf{\textit{Water insoluble}}} \\
\hline
Low MW lignin derived compound A & Dimethoxy stilbene & C\textsubscript{16}H\textsubscript{16}O\textsubscript{2} & 10.95 \\
\hline
Low MW lignin derived compound B & Dibenzofuran (representing diphenyl compounds) & C\textsubscript{12}H\textsubscript{8}O & 2.21 \\
\hline
Extractives & Dehydroabietic acid & C\textsubscript{20}H\textsubscript{30}O\textsubscript{2} & 2.99 \\
\hline
High MW lignin derived compound A & Oligomeric compounds with $\beta$-O-4 bond & C\textsubscript{20}H\textsubscript{26}O\textsubscript{8} & 9.15 \\
\hline
High MW lignin derived compound B & Phenylcoumaran compounds & C\textsubscript{21}H\textsubscript{26}O\textsubscript{8} & 1.99 \\
\hline
Nitrogen compounds & 2,4,6-trimethylpyridine & C\textsubscript{8}H\textsubscript{11}N & 0.070 \\
\hline
Sulfur compounds & Dibenzothiophene & C\textsubscript{12}H\textsubscript{8}S & 0.025 \\
\hline
\end{tabular}
\label{tab:model-compounds}
\end{table}

\begin{table}[htbp]
\centering
\caption{Boiling point and vapor pressure equation coefficients of user-defined compounds retrieved from \cite{jones2013process}.}
\begin{threeparttable}
% \scriptsize
\begin{tabular}{|l|l|l|p{4.5cm}|}
\hline
\textbf{Compound} & \textbf{Formula} & \textbf{Tbp, F} & \textbf{Coefficients for vapor pressure equation\tnote{1}} \\
\hline
Isoeugenol & C$_{10}$H$_{12}$O$_2$ & 507.2 & \begin{tabular}{l}
A = 50.75 \\
B = -9430.29 \\
C = -3.46 \\
D = -1.17E-6 \\
E = 2 \\
(Ref 2)
\end{tabular} \\
\hline
Levoglucosan & C$_6$H$_{10}$O$_5$ & 772.8 & \begin{tabular}{l}
A = 149.33 \\
B = -18972.6 \\
C = -16.826 \\
D = -4.57E-9 \\
E = 2 \\
(Ref 2)
\end{tabular} \\
\hline
Cellobiose & C$_{12}$H$_{22}$O$_{11}$ & 1234.2 & \begin{tabular}{l}
A = 53.01 \\
B = -19999.5 \\
C = -4.4294 \\
D = -1.5E-7 \\
E = 2
\end{tabular} \\
\hline
\begin{tabular}{l}
High MW lignin \\
derived compound A
\end{tabular} & C$_{20}$H$_{26}$O$_8$ & 1198.4 & \begin{tabular}{l}
A = 81.08 \\
B = -20000 \\
C = -7.036 \\
D = -2.279E-7 \\
E = 2
\end{tabular} \\
\hline
\begin{tabular}{l}
High MW lignin \\
derived compound B
\end{tabular} & C$_{21}$H$_{26}$O$_8$ & 1172.4 & \begin{tabular}{l}
A = 85.84 \\
B = -20000 \\
C = -7.72 \\
D = -3.84E-7 \\
E = 2
\end{tabular} \\
\hline
\end{tabular}
\begin{tablenotes}
\footnotesize
\item[1] Vp = exp[A + B/T + C ln(T) + D T$^E$], where Vp = vapor pressure (Pa), T = temperature (K)
\end{tablenotes}
\end{threeparttable}
\label{tab:app-user-define1}
\end{table}

\begin{table}[htbp]
\centering
\caption{Thermodynamic and critical properties of user-defined compounds retrieved from \cite{jones2013process}.}
% \scriptsize
\begin{tabular}{|l|l|p{1.1cm}|p{1.2cm}|p{1.1cm}|p{1.5cm}|p{1.5cm}|}
\hline
\textbf{Compound} & \textbf{Tc, F} & \textbf{Pc, psia} & \textbf{Vc, ft$^3$/ lbmol} & \textbf{Tm, F} & \textbf{$\Delta H_f$ (Btu/ lbmol)} & \textbf{$\Delta G_f$ (Btu/ lbmol)} \\
\hline
Isoeugenol  & 987.5 & 476.4 & 7.8 & 57.2 & -88516.7 & -16887.1 \\
\hline
Levoglucosan  & 1056 & 668.9 & 3.0 & 99.5 & -412740 & -255625 \\
\hline
Cellobiose  & 1281.5 & 340 & 6.8 & 462.2 & -958499 & -503968 \\
\hline
\begin{tabular}{l}
High MW \\ lignin derived \\ compound A
\end{tabular}  & 1482.2 & 241.9 & 15.9 & 218.3 & -498806 & -297590 \\
\hline
\begin{tabular}{l}
High MW \\ lignin derived \\ compound B
\end{tabular}  & 1468.2 & 243.7 & 16.0 & 252.3 & -496040 & -284159 \\
\hline
\end{tabular}
\label{tab:app-user-define2}
\end{table}

\FloatBarrier
\begin{longtable}{|p{0.4\textwidth}|p{0.5\textwidth}|}
\caption{Compound groups rising in hydroprocesses and their model compounds retrieved from \cite{jones2013process}.} \label{tab:compounds-from-hydroprocess}\\
\hline
\textbf{Groups} & \textbf{Model compounds} \\
\hline
\endfirsthead

\hline
\textbf{Generic names} & \textbf{Model compounds used to represent groups of similar compounds} \\
\hline
\endhead

1,2,3,4-Naphthalenes & 6-Hexyltetralin \\ \hline
Acids Intermediates & N-Butyric acid \\ \hline
Aromatic A & O-Xylene \\ \hline
Aromatic B & Benzene, 1-ethenyl-4-ethyl- \\ \hline
Ash & Ash \\ \hline
Bicyclos & 1,1-Bicyclohexyl \\ \hline
Carbon & Carbon \\ \hline
Carbonyl intermediates & 1,2-Propanediol \\ \hline
Extractive intermediates & HT Dehydroabietic acid \\ \hline
Cyclo C5 A & Ethyl-cyclopentane \\ \hline
Cyclo C5 B & 1-methyl-1-ethylcyclopentane \\ \hline
Cyclo C6 A & Cyclohexane \\ \hline
Cyclo C6 B & N-Butylcyclohexane \\ \hline
Cyclo C6-HCK product & Methyl-cyclohexane \\ \hline
Cyclo C7+ & 1,3-Dimethyladamantane \\ \hline
Diphenyl compounds & 1,2-Diphenylethane \\ \hline
Extractives & Dehydroabietic acid \\ \hline
High MW lignin compound intermediates & Indene \\ \hline
Indans & Indan \\ \hline
Indenes & 1H-Indene, 1,2,3-trimethyl- \\ \hline
Iso-paraffin A & 3-Methylhexane \\ \hline
Iso-paraffin B & 4-Methylnonane \\ \hline
Low MW lignin compound intermediates & Pinosylvin \\ \hline
Medium chain n-paraffins & N-Dodecane \\ \hline
Naphthalenes & 2,7-Dimethylnaphthalene \\ \hline
N-paraffin-HCK product A & N-Nonane \\ \hline
N-paraffin-HCK product B & N-Octane \\ \hline
Phenanthrenes & 4-Methylphenanthrene \\ \hline
Phenols & Thymol \\ \hline
Polynuclear aromatics & Naphthalene, 1-phenyl- \\ \hline
Pyrenes & Pyrene \\ \hline
Sand & Silicon Dioxide \\ \hline
Short chain n-paraffins & N-Hexane \\ \hline
Sugar intermediate A & Alpha-D-Glucose \\ \hline
Sugar intermediate B & Sorbitol \\ \hline
Sulfur agent & Dimethyl Sulfide \\
\hline
\end{longtable}

\newcommand{\COtwo}{CO\textsubscript{2}}
\newcommand{\COtwoe}{CO\textsubscript{2}e}
\newcommand{\Htwo}{H\textsubscript{2}}
\newcommand{\NHthree}{NH\textsubscript{3}}
\newcommand{\code}[1]{\texttt{#1}}

\section{Cost basis and financial parameters}\label{sec:S1}

\begin{itemize}
  \item Cost-year indices: \code{cost\_idx\_1Q2022 = 3230}, \code{cost\_idx\_Jones = 2250};
    Jones-era source costs escalated $\times(3230/2250)$ to 1Q-2022 USD. SMR uses its own index 2340.
  \item Financial: economic period 30 yr; IRR/discount 10\%; tax 35\%; salvage 0; straight-line
    depreciation; 8000 operating hours/yr ($\approx$91.3\% capacity factor).
  \item MJSP (annuity form):
    \begin{align*}
      \text{annuity} &= \frac{(1+\text{irr})^{-1}\,\bigl(1-(1+\text{irr})^{-n}\bigr)}{1-(1+\text{irr})^{-1}},\\
      \text{capital investment/yr} &= \frac{\text{CAPEX} - \text{annuity}\cdot\text{tax}\cdot\text{depreciation}}{\text{annuity}\,(1-\text{tax})},\\
      \text{MJSP} &= \frac{\text{min revenue} - \text{co-product sales}}{\text{annual SAF volume}}.
    \end{align*}
  \item SAF annual volume = stream 326SAF volumetric flow $\times$ 8000 h $\times$ 28.3168 L/cuft.
  \item Product/utility prices: 326SAF = 0 (the priced product); 316GAS (naphtha) 0.75757 \$/lb;
    336DIES (diesel) 0.6 \$/lb; electricity 19.05 \$/MMBtu; natural gas / fuel 14.38 \$/MMBtu.
\end{itemize}

%% ------------------------------------------------------------------
\section{Sub-unit cost interpolation}\label{sec:S2}

All five bolt-on sub-units use the same law
($r$ = stream/base stream): \textbf{CAPEX $\times r^{0.6}$} (six-tenths), \textbf{OPEX $\times r^{0.7}$},
stream and utility rates linear ($\times r$). Literature base cases are listed in Table~\ref{tab:S2-1}.
Hydrotreating catalyst charging: R-202/R-220 ruthenium LHSV 0.5,
price $60 \times 2.2$ \$/kg; R-225 molybdenum LHSV 0.22, price $15.5 \times 2.2$ \$/kg.

Its GHG emission is calculated from the simulated CCS cases and a constant ammonia-related emissions adjustment:
    \begin{equation}
    \mathrm{GHG}(a,c) =
    \mathrm{GHG}_{\mathrm{CCS}}(c)
    -a\left(\mathrm{GHG}_{\mathrm{grey}}-\mathrm{GHG}_{\mathrm{NH}_3}\right),
    \end{equation}
    where \(a\) denotes the green-ammonia substitution fraction and \(c\) denotes the CCS capture fraction. 

\begin{table}[htbp]
  \centering
  \caption{Literature base cases of the bolt-on sub-units.}
  \label{tab:S2-1}
  \resizebox{\textwidth}{!}{%
  \begin{tabular}{@{}llrrl@{}}
    \toprule
    Sub-unit & Reference & Base stream & Base CAPEX & Notes \\
    \midrule
    CCS (amine) & Rao \& Rubin 2002 & 862,853 lb/hr \COtwo & \$297.93 M & heat 3,315 MMBtu/hr; elec = 14\% of heat \\
    PSA \Htwo{} recovery & Byun 2020 & 700 m\textsuperscript{3}/hr & \$780,993 & OPEX \$5,078/yr; $\sim$68\% of OPEX = adsorbent/reactant \\
    R-510 hydrocracking & Jones 2013 & 9,600 lb/hr & \$5.27 M & driven by recycle stream 340 \\
    Off-gas reformer & Jones 2013 & (8846+995)/0.91 lb \Htwo & \$69.14 M (2011) & escalated $\times$3230/2250 \\
    SMR (+CCS) & Katebah 2022 & 45,922 & \$424 M / OPEX \$147 M & idx 2340; NG feed 156,136; \COtwo{} 55,107 \\
    \bottomrule
  \end{tabular}}
\end{table}


%% ------------------------------------------------------------------
\section{Scale-up and single-train limits}\label{sec:S3}

\begin{itemize}
  \item Scale set: [1, 2.5, 5, 7.5, 10, 20, 30]$\times$; main plant process 0.343~MMt/yr of lignocellulosic biomass.
  \item Pyrolyzer (R-110) single-train cap 700 MMBtu/hr = 10$\times$ base duty (70 MMBtu/hr at
    45.5 t/hr biomass); train count = ceil(base duty $\times$ scale / cap) $\rightarrow$ 1/1/1/1/1/2/3 across
    the scale set. Gas turbines, compressors, and expanders follow the same numbering-up.
  \item Ammonia-cracker reactor single-train cap 172 MMBtu/hr (distinct unit).
\end{itemize}

\begin{table}[htbp]
  \centering
  \caption{Per-unit single-train caps and replica counts. Number of parallel trains follows
    $n = \mathrm{ceil}(\text{total duty}/\text{cap})$; the cap equals the saturated per-replica load.}
  \label{tab:S3-1}
  \resizebox{\textwidth}{!}{%
  \begin{tabular}{@{}lllll@{}}
    \toprule
    Unit & Service & Sizing metric & Single-train cap & Replicas [1, 2.5, 5, 7.5, 10, 20, 30]$\times$ \\
    \midrule
    C-101 & Compressor & Volumetric flow & 15.81$\times$10\textsuperscript{6} cuft/hr (2772 HP) & 1 / 1 / 1 / 1 / 1 / 2 / 3 \\
    C-110 & Compressor & Volumetric flow & 17.57$\times$10\textsuperscript{6} cuft/hr & 1 / 1 / 1 / 2 / 2 / 3 / 5 \\
    C-147 & Compressor & Volumetric flow & 17.77$\times$10\textsuperscript{6} cuft/hr & 1 / 1 / 1 / 2 / 2 / 4 / 6 \\
    C-401 & Compressor & Volumetric flow & 15.23$\times$10\textsuperscript{6} cuft/hr (192,935 HP) & 1 / 1 / 1 / 1 / 1 / 2 / 3 \\
    MC-610 & Compressor & Volumetric flow & 17.38$\times$10\textsuperscript{6} cuft/hr & 1 / 2 / 3 / 4 / 5 / 10 / 20 $\dagger$ \\
    C-231 & Expander & Volumetric flow & 24,981 cuft/hr (1944 HP) & 2 / 6 / 10 / 14 / 18 / 36 / 54 $\dagger$ \\
    R-110 & Combustion chamber (pyrolyzer) & Heat absorbed & 447.5 MMBtu/hr (base 44.75) & 1 / 1 / 1 / 1 / 1 / 2 / 3 \\
    E-616 & Heat exchanger & Heat duty & 955.4 MMBtu/hr & 1 / 1 / 1 / 1 / 1 / 2 / 3 \\
    GT-612 & Gas turbine & Power output & 441,073 hp & 1 / 1 / 2 / 2 / 2 / 4 / 6 \\
    \bottomrule
  \end{tabular}}
\end{table}

\begin{itemize}
  \item The R-110 cap of 447.5 MMBtu/hr (10$\times$ the 44.75 MMBtu/hr base) supersedes the
    700 MMBtu/hr figure listed above; $\mathrm{ceil}(\text{scale} \times 44.75 / 447.5)$ reproduces its
    1/1/1/1/1/2/3 pattern exactly.
  \item $\dagger$ MC-610 and C-231 are \textbf{over-provisioned} relative to the strict
    $\mathrm{ceil}(\text{total}/\text{cap})$ rule: MC-610 carries 20 trains at 30$\times$ where the rule gives 15
    (total/cap = 15.0 exactly), and C-231 carries 6 and 10 trains at 2.5$\times$ and 5$\times$ where the rule
    gives 5 and 9. Every other unit matches the rule exactly.
\end{itemize}

%% ------------------------------------------------------------------
\section{Ammonia cracker detail}\label{sec:S4}

\begin{figure}[htbp]
    \centering
    \includegraphics[width=\textwidth]{BiomassLCA-Ammonia (2).png}
    \caption{Simplified process flow diagram of the ammonia cracking unit: green, blue, and orange line refer to the streams with ammonia feed, hydrogen-rich product, and residual gas, respectively.}
    \label{fig:pfd-ammonia}
\end{figure}

Ammonia cracker is simplified in Fig.~\ref{fig:pfd-ammonia}.
\begin{itemize}
  \item Cracker scale index $n$ = main \Htwo{} demand / 6570.40 lb/hr; non-grid $n$ linearly
    interpolated between bracketing APEA runs.
  \item Utilities = electricity; raw material = ammonia; fertilizer credited from the Product
    Summary sheet. Conversion / \Htwo-recovery ratios live inside the Aspen model (APEA export
    gives \NHthree-in vs \Htwo-out); GHSV/kinetics from \cite{cho2024ammonia}.
  \item Grey-\Htwo{} benchmark 0.68 \$/lb \Htwo. SMR uses natural-gas price 0.092306 \$/lb (US), escalated.
  \item 
\end{itemize}

%% ------------------------------------------------------------------
\section{CCS detail}\label{sec:S5}

\begin{itemize}
  \item Sequestered mass $\text{co2}_\text{seq} = \text{co2}_\text{left}(\text{623CO2})\cdot\text{ccs portion} + \text{co2}_\text{EGR}(\text{622X})$;
    622X captured at 100\%, 623CO2 at ccs portion.
  \item Capture OPEX: Fuel $= \max(\text{heat} - \text{plant heat gen}, 0)\cdot 14.38\cdot 8000$;
    Electricity $= \text{elec}\cdot 19.05\cdot 8000$; T\&S $= \text{co2}_\text{seq}/2204.62\cdot 8000\cdot(\text{co2 TS cost per ton}$).
  \item \COtwo{} T\&S price: default 0 (off) in base runs; literature range 4--45 \$/t; manuscript basis
    10 \$/t. 
  \item ccs portion sweeps: \{0, 0.25, 0.5, 0.75, 1.0\}; 0--0.9 in 0.1 steps; scale-up scenario fixed at 0.60.
\end{itemize}


%% ------------------------------------------------------------------
% \section{LCA inventories, allocation, and SAF criterion}\label{sec:S7}

% In this paper, LCA on the environmental impact of the developed process designs was conducted following the CORSIA methodology \cite{ICAO2023CORSIA} as illustrated in Fig. \ref{fig:lca-scope}.
% The International Civil Aviation Organization (ICAO) has proposed the method as a standard so that the SAF industry can consistently and transparently monitor its GHG emission intensities. 
% % It was formerly proposed by the International Civil Aviation Organization (ICAO) as a standard to guide the SAF industry to monitor their GHG emission intensities consistently and transparently. 
% CORSIA life cycle emission value is the sum of the direct emissions along the supply chains of individual SAFs and the estimated emissions for induced land use change (ILUC) \cite{prussi2021corsia}.
% It is also referred to as WTWa emissions since its scope spans from feedstock preparation to fuel combustion on the aircraft. 



%% ------------------------------------------------------------------
\section{Biomass price model in supply-availability context}\label{sec:S7}

Domestic feedstock ceilings used to convert plant scale to a share of aviation-fuel demand: forest residue 3,127,378 t/yr; log (2,843,071 + 3,274,200) t/yr. 

\textbf{FR Price (reference).} A constant forest-residue price of 107.3~\$/t, independent of scale; this is the optimistic floor.

\textbf{FR--Log Linear} (naive Galik supply curve; \citealp{galik}). The forest-residue price holds until regional forest-residue supply (2.84~MMt/yr) is exhausted, beyond which the delivered price rises linearly with demand, price = supply $\times$ 25.08 + 33.30~\$/t (supply in MMt/yr), as logging residue is drawn in.

\textbf{Fixed Elasticity} (Galik supply curve, supply elasticity 0.3). Follows the FR--Log Linear curve until the combined forest-residue-plus-log supply (6.12~MMt/yr) is exhausted, then extrapolates linearly under a fixed supply elasticity of 0.3 (log reference price 192.4~\$/t), giving the steepest scale penalty.

\textbf{Scale-up evaluation.} MJSP is evaluated at seven plant scales: 0.041, 0.104, 0.207, 0.311, 0.414, 0.828 and 1.243~MMt/yr SAF. At the forest-residue floor price, the base-case (grey-\ce{H2}) MJSP at each scale is broken down into capital investment, raw materials, utilities, fixed operating costs (operating labor, maintenance, operating charges, plant overhead, and general and administrative costs) and co-product sales, which enter as a negative term (Table~\ref{tab:S9-1}). Each of the three supply-curve models then sets the delivered biomass price at each scale, and MJSP is recomputed for three configurations: the grey-\ce{H2} base, full green ammonia, and full CCS (Tables~\ref{tab:S9-2} and \ref{tab:S9-3}). Biomass is costed identically in all three configurations, so a given price model shifts each of them by the same biomass-cost increment. The minimum-MJSP scale is read where the falling conversion cost and the rising feedstock cost cross.
%% ------------------------------------------------------------------
%% ------------------------------------------------------------------
\section{Data tables for the Results figures}\label{sec:S9}

Numeric values behind the main-text Results figures.
MJSP in \$/L SAF (1Q-2022 USD); GHG in g \COtwoe/MJ (WTWa GHG-100, Allocation, Korea grid mix).

\begin{table}[htbp]
  \centering
  \caption{Base-case (grey \Htwo, flat forest-residue price) MJSP breakdown vs SAF production
    (base-case breakdown figure).
    Components in \$/L SAF; MJSP = sum of components.}
  \label{tab:S9-1}
  \resizebox{\textwidth}{!}{%
  \begin{tabular}{@{}rrrrrrrrrrr@{}}
    \toprule
    \shortstack{SAF production\\{[MMt/yr]}} & \shortstack{Capital\\investment} & \shortstack{Raw\\materials} & Utilities &
    \shortstack{Operating\\labor} & Maintenance & \shortstack{Operating\\charges} & \shortstack{Plant\\overhead} & G\&A &
    \shortstack{Co-product\\sales} & MJSP \\
    \midrule
    0.041 & 1.152 & 1.284 & 0.676 & 0.036 & 0.110 & 0.009 & 0.073 & 0.172 & $-1.860$ & 1.651 \\
    0.104 & 0.936 & 1.269 & 0.603 & 0.014 & 0.093 & 0.004 & 0.054 & 0.160 & $-1.860$ & 1.271 \\
    0.207 & 0.857 & 1.263 & 0.480 & 0.007 & 0.086 & 0.002 & 0.046 & 0.147 & $-1.860$ & 1.028 \\
    0.311 & 0.801 & 1.262 & 0.419 & 0.005 & 0.083 & 0.001 & 0.044 & 0.142 & $-1.860$ & 0.897 \\
    0.414 & 0.750 & 1.261 & 0.381 & 0.004 & 0.078 & 0.001 & 0.041 & 0.138 & $-1.860$ & 0.793 \\
    0.828 & 0.719 & 1.260 & 0.340 & 0.003 & 0.077 & 0.001 & 0.040 & 0.134 & $-1.860$ & 0.714 \\
    1.243 & 0.700 & 1.259 & 0.324 & 0.002 & 0.079 & 0.000 & 0.040 & 0.133 & $-1.860$ & 0.676 \\
    \bottomrule
  \end{tabular}}
\end{table}

\begin{table}[htbp]
  \centering
  \caption{Biomass price and grey-\Htwo{} base MJSP under three biomass price models
    (price-model figure). MJSP in \$/L SAF.}
  \label{tab:S9-2}
  \resizebox{\textwidth}{!}{%
  \begin{tabular}{@{}rrrrrrr@{}}
    \toprule
    & \multicolumn{3}{c}{Biomass price [\$/lb]} & \multicolumn{3}{c}{MJSP grey \Htwo{} [\$/L]} \\
    \cmidrule(lr){2-4}\cmidrule(l){5-7}
    SAF production [MMt/yr] & FR Price & FR-Log Linear & Fixed Elasticity & FR Price & FR-Log Linear & Fixed Elasticity \\
    \midrule
    0.041 & 0.0486 & 0.0486 & 0.0486 & 1.651 & 1.651 & 1.651 \\
    0.104 & 0.0486 & 0.0486 & 0.0486 & 1.271 & 1.271 & 1.271 \\
    0.207 & 0.0486 & 0.0486 & 0.0486 & 1.028 & 1.028 & 1.028 \\
    0.311 & 0.0486 & 0.0486 & 0.0486 & 0.897 & 0.897 & 0.897 \\
    0.414 & 0.0486 & 0.0527 & 0.0527 & 0.793 & 0.849 & 0.849 \\
    0.828 & 0.0486 & 0.0903 & 0.1064 & 0.714 & 1.290 & 1.515 \\
    1.243 & 0.0486 & 0.1278 & 0.2360 & 0.676 & 1.773 & 3.272 \\
    \bottomrule
  \end{tabular}}
\end{table}

\begin{table}[htbp]
  \centering
  \caption{MJSP of the full green-\NHthree{} and full-CCS configurations under three biomass price models
    (price-model figure; full green \NHthree{} = 100\% green-\NHthree{} hydrogen,
    full CCS = CCS portion 1.0). MJSP in \$/L SAF.}
  \label{tab:S9-3}
  \resizebox{\textwidth}{!}{%
  \begin{tabular}{@{}rrrrrrr@{}}
    \toprule
    & \multicolumn{3}{c}{MJSP full green \NHthree{} [\$/L]} & \multicolumn{3}{c}{MJSP full CCS [\$/L]} \\
    \cmidrule(lr){2-4}\cmidrule(l){5-7}
    SAF production [MMt/yr] & FR Price & FR-Log Linear & Fixed Elasticity & FR Price & FR-Log Linear & Fixed Elasticity \\
    \midrule
    0.041 & 3.495 & 3.495 & 3.495 & 2.275 & 2.275 & 2.275 \\
    0.104 & 3.062 & 3.062 & 3.062 & 1.852 & 1.852 & 1.852 \\
    0.207 & 2.802 & 2.802 & 2.802 & 1.586 & 1.586 & 1.586 \\
    0.311 & 2.661 & 2.661 & 2.661 & 1.445 & 1.445 & 1.445 \\
    0.414 & 2.544 & 2.600 & 2.600 & 1.334 & 1.390 & 1.390 \\
    0.828 & 2.443 & 3.020 & 3.244 & 1.242 & 1.818 & 2.043 \\
    1.243 & 2.401 & 3.498 & 4.997 & 1.198 & 2.295 & 3.794 \\
    \bottomrule
  \end{tabular}}
\end{table}

\begin{table}[htbp]
  \centering
  \caption{MJSP, GHG, Pareto status, sensitivity band, and abatement cost of the 30 green-\NHthree{} share
    $\times$ CCS portion configurations at 0.414 MMt/yr SAF.  Abatement cost relative to CCS 60\%
    (75.14 g \COtwoe/MJ, \$0.977/L); -- where GHG is not below CCS 60\%.}
  \label{tab:S9-4}
  \small
  \begin{tabular}{@{}rrrrcr@{}}
    \toprule
    \shortstack{Green \NHthree\\share [\%]} & \shortstack{CCS\\portion [\%]} & \shortstack{MJSP\\{[\$/L]}} &
    \shortstack{GHG\\{[g \COtwoe/MJ]}} & Pareto &
    \shortstack{Abatement cost\\{[\$/t \COtwo]}} \\
    \midrule
    0 & 0 & 0.792 & 100.83 & yes & -- \\
    0 & 25 & 0.855 & 89.93 & yes & -- \\
    0 & 50 & 0.913 & 78.99 & yes & -- \\
    0 & 75 & 1.111 & 69.36 & yes & 733 \\
    0 & 100 & 1.334 & 59.96 & yes & 742 \\
    10 & 0 & 0.977 & 97.02 & no & -- \\
    10 & 25 & 1.039 & 86.12 & no & -- \\
    10 & 50 & 1.098 & 75.18 & yes & -- \\
    10 & 75 & 1.296 & 65.55 & yes & 1,049 \\
    10 & 100 & 1.518 & 56.15 & yes & 900 \\
    25 & 0 & 1.240 & 91.31 & no & -- \\
    25 & 25 & 1.302 & 80.40 & no & -- \\
    25 & 50 & 1.361 & 69.46 & no & 2,138 \\
    25 & 75 & 1.559 & 59.83 & no & 1,201 \\
    25 & 100 & 1.782 & 50.44 & yes & 1,029 \\
    50 & 0 & 1.680 & 81.78 & no & -- \\
    50 & 25 & 1.742 & 70.88 & no & 5,671 \\
    50 & 50 & 1.800 & 59.94 & no & 1,711 \\
    50 & 75 & 1.998 & 50.31 & yes & 1,299 \\
    50 & 100 & 2.221 & 40.91 & yes & 1,148 \\
    75 & 0 & 2.116 & 72.26 & no & 12,511 \\
    75 & 25 & 2.178 & 61.35 & no & 2,752 \\
    75 & 50 & 2.236 & 50.41 & no & 1,609 \\
    75 & 75 & 2.434 & 40.78 & yes & 1,340 \\
    75 & 100 & 2.657 & 31.39 & yes & 1,213 \\
    100 & 0 & 2.544 & 62.74 & no & 3,991 \\
    100 & 25 & 2.606 & 51.83 & no & 2,207 \\
    100 & 50 & 2.665 & 40.89 & no & 1,556 \\
    100 & 75 & 2.862 & 31.26 & yes & 1,357 \\
    100 & 100 & 3.085 & 21.86 & yes & 1,250 \\
    \bottomrule
  \end{tabular}
\end{table}

\begin{table}[htbp]
  \centering
  \caption{Grey-\Htwo{} price at which Green \NHthree{} 10\% + CCS 50\% matches CCS 60\%.}
  \label{tab:S9-5}
  \small
  \begin{tabular}{@{}rrrc@{}}
    \toprule
    Green \NHthree{} price [\$/t] & Crossover grey-\Htwo{} price [\$/kg] & MJSP at crossover [\$/L] & Inside 1--4 \$/kg range \\
    \midrule
    270 & 1.59 & 1.016 & yes \\
    441 & 2.91 & 1.568 & yes \\
    635 (base) & 4.40 & 2.195 & no \\
    1,000 & 7.21 & 3.373 & no \\
    \bottomrule
  \end{tabular}
\end{table}

\begin{table}[htbp]
  \centering
  \caption{Mean $|$SHAP$|$ of MJSP drivers at the two cheapest compliant points in \$/L SAF: 128 Latin-hypercube TEA runs per point,
    gradient-boosted surrogate ($R^2 > 0.999$); -- = driver not sampled.}
  \label{tab:S9-6}
  \small
  \begin{tabular}{@{}lrr@{}}
    \toprule
    Driver & CCS 60\% & Green \NHthree{} 10\% + CCS 50\% \\
    \midrule
    Grey \Htwo & 0.310 & 0.274 \\
    CAPEX & 0.103 & 0.102 \\
    Co-product & 0.097 & 0.083 \\
    Electricity & 0.052 & 0.055 \\
    Green \NHthree & -- & 0.031 \\
    \COtwo{} T\&S & 0.025 & 0.023 \\
    Capacity & 0.019 & 0.028 \\
    Biomass & 0.003 & 0.002 \\
    \bottomrule
  \end{tabular}
\end{table}

\begin{table}[htbp]
  \centering
  \caption{One-at-a-time MJSP spans at the two cheapest compliant points in \$/L SAF; input ranges as in the tornado-driver analysis.}
  \label{tab:S9-7}
  \small
  \begin{tabular}{@{}llrrrr@{}}
    \toprule
    Point & Driver & Base MJSP & MJSP at low input & MJSP at high input & Span \\
    \midrule
    CCS 60\% & Grey \Htwo & 0.977 & 0.767 & 2.026 & 1.258 \\
    CCS 60\% & Co-product & 0.977 & 1.224 & 0.730 & 0.494 \\
    CCS 60\% & CAPEX & 0.977 & 0.741 & 1.213 & 0.472 \\
    CCS 60\% & Electricity & 0.977 & 1.134 & 0.820 & 0.313 \\
    CCS 60\% & Capacity & 0.977 & 1.174 & 0.987 & 0.188 \\
    CCS 60\% & \COtwo{} T\&S & 0.977 & 0.950 & 1.137 & 0.187 \\
    CCS 60\% & Biomass & 0.977 & 0.977 & 1.029 & 0.052 \\
    \midrule
    Green \NHthree{} 10\% + CCS 50\% & Grey \Htwo & 1.098 & 0.909 & 2.042 & 1.133 \\
    Green \NHthree{} 10\% + CCS 50\% & Co-product & 1.098 & 1.344 & 0.851 & 0.494 \\
    Green \NHthree{} 10\% + CCS 50\% & CAPEX & 1.098 & 0.858 & 1.337 & 0.480 \\
    Green \NHthree{} 10\% + CCS 50\% & Electricity & 1.098 & 1.258 & 0.937 & 0.321 \\
    Green \NHthree{} 10\% + CCS 50\% & Green \NHthree & 1.098 & 0.980 & 1.215 & 0.234 \\
    Green \NHthree{} 10\% + CCS 50\% & Capacity & 1.098 & 1.313 & 1.108 & 0.205 \\
    Green \NHthree{} 10\% + CCS 50\% & \COtwo{} T\&S & 1.098 & 1.073 & 1.238 & 0.165 \\
    Green \NHthree{} 10\% + CCS 50\% & Biomass & 1.098 & 1.149 & 1.149 & 0.052 \\
    \bottomrule
  \end{tabular}
\end{table}

\bibliographystyle{cas-model2-names}
\bibliography{main-sc}

\end{document}



% To print the credit authorship contribution details
\printcredits

\section*{Acknowledgements}
This work was supported by the Settlement Research Fund of UNIST (Ulsan National Institute of Science \& Technology).

%% Loading bibliography style file
%\bibliographystyle{model1-num-names}
\bibliographystyle{cas-model2-names}

% Loading bibliography database
\bibliography{main-sc}

% Biography
%\bio{}
% Here goes the biography details.
%\endbio

%\bio{pic1} 
% Here goes the biography details.
%\endbio

\end{document}

